% =====================================================================
%  03_estado-del-arte.tex — Trabajos y soluciones relacionadas
%  Versión preliminar: completar con más fuentes antes del cierre.
% =====================================================================
\chapter{Estado del arte}
\label{cap:estado}

El problema de \saru{} se relaciona con el \textbf{reconocimiento
automático de matrículas} (ALPR, por sus siglas en inglés), un campo con
amplia literatura. La diferencia principal es que \saru{} no lee una placa
con formato estándar, sino un número económico pintado en distintas zonas
de la carrocería, y que además debe decidir la dirección del movimiento.

\section{Trabajos relacionados}

\begin{description}[style=nextline]
  \item[Revisión de métodos de ALPR] Shashirangana et~al. revisan los
        métodos de reconocimiento de matrículas y los organizan por etapa:
        detección de la placa, segmentación y reconocimiento de caracteres.
        Señalan como retos la variación de iluminación, el ángulo de la
        cámara y la oclusión parcial~\cite{shashirangana2021alpr}.
  \item[ALPR en tiempo real con YOLO] Laroca et~al. proponen un sistema de
        reconocimiento de matrículas basado por completo en detectores
        YOLO, con etapas en cascada para el vehículo, la placa y los
        caracteres, y muestran que puede operar en tiempo
        real~\cite{laroca2018yolo}.
\end{description}

\section{Comparación con \saru{}}

\begin{longtable}{@{}L{3.6cm}L{3.3cm}L{3.3cm}L{3.6cm}@{}}
  \caption{Comparación entre los trabajos relacionados y \saru{}.}
  \label{tab:comparacion} \\
  \toprule
  \textbf{Característica} & \textbf{ALPR general~\cite{shashirangana2021alpr}}
    & \textbf{YOLO en cascada~\cite{laroca2018yolo}} & \textbf{\saru{}} \\
  \midrule
  \endfirsthead
  \toprule
  \textbf{Característica} & \textbf{ALPR general} & \textbf{YOLO en cascada}
    & \textbf{\saru{}} \\
  \midrule
  \endhead
  Objeto que se lee & Matrícula & Matrícula & Número económico pintado \\
  Zonas de lectura & Una (placa) & Una (placa) & Delantera, lateral o
    trasera \\
  Dirección del movimiento & No aplica & No & Sí (línea virtual) \\
  Procesamiento & Variable & Por verificar & En el borde (Jetson) \\
  Evidencia y avisos & No aplica & No & Fotos, video y avisos \\
  \bottomrule
\end{longtable}

\begin{nota}[title={Por completar}]
Esta sección es preliminar. Falta agregar soluciones comerciales de control
de acceso vehicular, trabajos sobre identificación de autobuses de
transporte público y conteo de vehículos con líneas virtuales, y revisar
que la descripción de cada fuente coincida con el artículo original.
\end{nota}
